\documentclass[10pt]{article}
\usepackage[margin=0.7in,top=0.55in,bottom=0.55in]{geometry}
\usepackage{amsmath,amssymb,booktabs,microtype,xcolor,url}
\usepackage[colorlinks=true,linkcolor=blue!60!black,urlcolor=blue!60!black,pdftitle={FRE-GY-9743 Assignment 8},pdfauthor={Raj Pawar}]{hyperref}
\newcommand{\code}[1]{\texttt{#1}}
\setlength{\parindent}{0pt}\setlength{\parskip}{3pt}\emergencystretch=3em\pagestyle{plain}
\makeatletter\g@addto@macro\normalsize{\setlength{\abovedisplayskip}{5pt}\setlength{\belowdisplayskip}{5pt}\setlength{\abovedisplayshortskip}{3pt}\setlength{\belowdisplayshortskip}{3pt}}\makeatother
\makeatletter\renewcommand{\section}{\@startsection{section}{1}{\z@}{8pt plus 2pt}{3pt}{\normalfont\large\bfseries}}\makeatother
\begin{document}
{\Large\bfseries FRE-GY-9743 Assignment 8: SOFR OIS Calibration and Hedging Risk}\\[1pt]
Raj Pawar, NYU Tandon School of Engineering, November 21, 2026.\\
Code: \href{https://github.com/pawarraj8888/FRE-GY-9743-Assignments-8/tree/Raj_Pawar_Branch}{github.com/pawarraj8888/FRE-GY-9743-Assignments-8}, branch \code{Raj\_Pawar\_Branch}.\\
Project page: \href{https://pawarraj8888.github.io/FRE-GY-9743-Assignments-8/}{pawarraj8888.github.io/FRE-GY-9743-Assignments-8}.
\vspace{2pt}\hrule

\section*{1. What the notebook does}
The workbook gives 17 SOFR OIS par quotes, 1Y to 50Y, on the valuation date 2026-10-09 (447.8 bp at 1Y rising to 501.3 bp at 20Y and
back to 466.5 bp at 50Y). The SOFR curve is an instantaneous forward rate curve, piecewise constant and left-continuous with flat
extrapolation, with one node per quote placed at that swap's last payment date; the funding curve is SOFR plus a zero spread with a
single node held at zero. The library bootstraps: components are built in order, and for each OIS in maturity order the node forward is
solved with Brent's method so that the swap has zero PV. Because the forward is piecewise constant, a new node only changes the forward
on its own segment, so each solve is one-dimensional and leaves the earlier instruments repriced. After the five pieces below were
filled in, all 17 quotes reprice to within 5.3e-10 bp.

\section*{2. The five pieces}
\textbf{2a, 2b: annuity and par rate} (\code{yield\_curve/valuation\_engine.py}, \code{ValuationEngineProductRfrSwap.calculate\_value}). The fixed leg pays $NK\alpha_i$ on each
remaining coupon date $p_i$, so its PV already calculated by the leg engine is $V_{\mathrm{fix}}=NK\sum_i\alpha_iD(p_i)$ and the annuity
per unit notional is $A=V_{\mathrm{fix}}/(NK)$; the coupon $K$ is the market quote for calibration instruments and nonzero in every
use here. The par rate is the coupon that makes the swap worth zero, $R=V_{\mathrm{float}}/(NA)$. For the 9Y6M swap of Section 3,
$A=7.5863$ and $R=4.8273\%$, so receiving $4.5\%$ is worth $(K-R)\,NA=$ USD -24,829.03, the number the notebook prints.

\textbf{2c: par-rate gradient} (same file, \code{grad\_at\_par}). With $K$ held fixed, $R=V_{\mathrm{float}}/(NA)$ and $NA=V_{\mathrm{fix}}/K$, so
by the quotient rule
\[
\nabla_{X^I}R=\frac{\nabla_{X^I}V_{\mathrm{float}}}{NA}-R\,\frac{\nabla_{X^I}V_{\mathrm{fix}}}{V_{\mathrm{fix}}},
\]
where $X^I$ is the vector of internal state variables (the node forwards of every component, flattened in the model's component
order). The two leg engines already know how to accumulate their PV gradients block by block, so the implementation calls each of
them once with the matching scalar, which keeps the blocks aligned with the state layout.

\textbf{3: calibration Jacobian} (\code{yield\_curve/yield\_curve\_model.py}, \code{YieldCurve.calculate\_model\_jacobian}). $J=\partial X^M/\partial X^I$ has one row per
calibration instrument: the OIS rows are the gradients from 2c, the funding row comes from the zero-spread engine, and each row is the
flattened block list. Here $J$ is $18\times18$ (17 OIS nodes and one funding node). Its OIS block is lower triangular, because quote
$i$ only depends on nodes up to its own maturity, with a positive diagonal: raising a swap's last-segment forward raises its par rate.

\textbf{4: quote risk} (\code{valuation/utilities.py}, \code{risk\_calculation}). The engines give $\nabla_{X^I}V$. The market quotes are a function of the state,
$X^M=F(X^I)$ with $J=\partial F/\partial X^I$ invertible, so a move $dX^M$ in quotes is produced by $dX^I=J^{-1}dX^M$ and the chain rule
gives
\[
\nabla_{X^M}V=\nabla_{X^I}V\,J^{-1},
\]
which I compute by solving $J^{\top}r^{\top}=(\nabla_{X^I}V)^{\top}$ rather than forming the inverse. The values stay in decimal-rate
units and the report multiplies by $10^{-4}$ to show USD per basis point.

\textbf{Hedge.} With $B_{ij}$ the quote risk of one million receive-fixed par OIS $i$ to quote $j$ (a $17\times17$ matrix) and $r$ the
target's OIS quote risk, the weights solve $B^{\top}w=-r^{\top}$; $w_i$ is the receive-fixed notional in millions, negative meaning pay
fixed. Since $B$ is square and nonsingular the residual $r+w^{\top}B$ is zero to machine precision.

\section*{3. Results and how to read them}
\begin{center}\footnotesize
\begin{tabular}{@{}lrrrr@{}}\toprule
OIS & Quote (bp) & Target risk (USD/bp) & Par OIS own-quote risk (USD/bp) & Hedge notional (USD, receive $>0$) \\ \midrule
1Y & 447.795 & 1.336 & -97.77 & 13,666 \\
2Y & 465.672 & 0.382 & -189.91 & 2,011 \\
3Y & 470.730 & 0.597 & -277.73 & 2,151 \\
4Y & 472.032 & 1.038 & -361.78 & 2,870 \\
5Y & 473.200 & 0.823 & -441.91 & 1,863 \\
6Y & 474.988 & 1.766 & -519.13 & 3,403 \\
7Y & 477.000 & 1.567 & -591.67 & 2,648 \\
8Y & 479.203 & 3.531 & -660.73 & 5,345 \\
9Y & 481.600 & -355.339 & -726.42 & -489,162 \\
10Y & 484.300 & -401.432 & -789.18 & -508,671 \\
12Y & 489.725 & 0.000 & -905.49 & 0 \\
15Y & 496.507 & 0.000 & -1,057.51 & 0 \\
20Y & 501.255 & 0.000 & -1,264.03 & 0 \\
25Y & 498.939 & 0.000 & -1,425.45 & 0 \\
30Y & 493.568 & 0.000 & -1,554.29 & 0 \\
40Y & 480.540 & 0.000 & -1,747.51 & 0 \\
50Y & 466.540 & 0.000 & -1,886.18 & 0 \\
\bottomrule\end{tabular}\end{center}

\textbf{Why the risk sits at 9Y and 10Y.} A receive-fixed swap loses when rates rise, so its net risk is negative, and it is
concentrated in the two quotes whose forward segments carry the swap's last cashflows: -355.34 USD/bp at 9Y and -401.43 at 10Y; the
earlier quotes carry small positive residuals explained below.
The bootstrap explains the shape. Raising the 9Y quote by one basis point with every other quote fixed lifts the forward on the
8Y--9Y segment by 10.8 bp (the 9Y par rate is a discount-weighted average of nine yearly forwards, so the last year has to carry
roughly eleven times the move) and, since the 10Y swap must still
reprice, lowers the 9Y--10Y segment by 11.3 bp. The 9Y6M swap pays floating on all of the first segment and on half of the
second, so the two moves do not cancel and the swap loses about $10^6\times(\text{10.8}-\tfrac12\times\text{11.3})\,\text{bp}
\times 10^{-4}\approx 510$ USD before discounting, the right order for the -355.34 printed (the exact number also carries the change
in the funding discount factors). Raising the 10Y quote lifts the 9Y--10Y segment by 12.3 bp
and lowers the 10Y--12Y segment, which the swap never touches, hence -401.43. For a quote inside the swap's life, say 5Y, the
segment it lifts (5.4 bp on 4Y--5Y) and the segment it lowers (5.7 bp on 5Y--6Y) are both fully paid by the swap, so the two
effects cancel to first order and only a small positive residual from discounting is left (0.82 USD/bp); the quotes beyond 10Y
cannot move the swap at all, so their risk is exactly zero.

\textbf{The hedge.} The solver returns pay-fixed positions of USD 489,162 at 9Y and USD 508,671 at 10Y, together close to one million
(997,833 in total, 49\% of it at 9Y, as a 9Y6M maturity sits about half way between the two nodes), plus small receive-fixed
positions of a few thousand at 1Y to 8Y that absorb the residual risks, and nothing beyond 10Y. The own-quote risk of a one million
par OIS equals $-NA\times10^{-4}$, its PV01, which is why the diagonal of $B$ grows with maturity (-97.77 at 1Y, -1,886.18 at 50Y). The
hedge is first order: when I move the 9Y quote by a full basis point and recalibrate, the naked swap moves by USD -355.42 and the
hedged book by USD -0.12 (for 10Y: -401.52 and -0.12), the leftover being convexity.

\section*{4. Checks}
\code{tests/test\_ois\_calibration\_and\_risk.py} has 20 tests: every calibration swap reprices to its quote and has zero PV;
nodes at the last payment dates in maturity order; the annuity equals the hand sum of accrual times discount factor from the cashflow
report; the par rate equals the floating PV over the annuity, prices the swap to zero, and does not depend on direction or notional;
the reference PV; the par-rate gradient against central differences in the state and against the quotient rule of the two leg
gradients; the Jacobian's shape, triangular structure and rows against finite differences of the quotes; the quote risk against the
shipped reference, against the gradient times the explicit inverse, and against bump-and-recalibrate of the 1Y, 8Y, 9Y, 10Y and 12Y
quotes (and, for a forward-starting 2Yx5Y, a 15Y and a negative-notional swap, the gradient and the 2Y, 7Y, 12Y and 15Y risks); the
funding-spread row of the report, which is the swap's sensitivity to the zero spread between the funding and SOFR curves; signs,
locality and the pay/receive mirror; a zero-coupon swap's annuity and par rate; a par OIS's own-quote risk as its PV01 with zero risk elsewhere; the report labels;
and the hedge (reference notionals, zero residual, the 9Y/10Y split, and a real one basis point move of the hedged book). The notebook
is the professor's, executed top to bottom, and its tables equal the reference outputs he shipped in it to the printed precision
(the $10^{-10}$ bp repricing errors differ in their last digits).
\end{document}
